\documentclass[12pt]{article}
\usepackage[utf8]{inputenc}
\usepackage[T1]{fontenc}
\usepackage{hyperref}
\title{Yabula Yabula Language Resources for K-12 Curriculum Development}
\author{Compiled by Claude Sonnet 5.5 (claude-sonnet-5-5) for Andr\'as Kornai}
\date{}
\begin{document}
\maketitle

\section{Introduction}
Yabula Yabula (ISO 639-3 yxy; Glottolog yabu1234; AIATSIS S38) is an extinct Pama-Nyungan language of the Murray-Goulburn area of Victoria and New South Wales. Its classification is disputed. Dixon treated it as an isolate, and Glottolog treats it as a dialect of Yorta Yorta. Bowe and Morey (1999) report that it shares only about 44\% of its vocabulary with Yorta Yorta. Documentation is thin.

\section{Primary Educational Resources}
No Yabula Yabula primers or school programmes were found. The Victorian Curriculum has an Aboriginal Languages framework, and the VCAA document surfaced in search results. I did not verify its details.

\section{Linguistic and Reference Materials}
\begin{itemize}
\item Bowe, Heather and Morey, Stephen (1999), \emph{The Yorta Yorta (Bangerang) Language of the Murray Goulburn including Yabula Yabula}, Pacific Linguistics C-154. The search summary says it is open access under CC BY-SA 4.0 from the ANU repository. This is the main reference. \url{https://openresearch-repository.anu.edu.au/server/api/core/bitstreams/c0ddf545-58b8-46a9-a087-fb46f254b7d5/content}
\item Wikipedia, ``Yabula-Yabula language''. \url{https://en.wikipedia.org/wiki/Yabula-Yabula_language}
\item Museums Victoria, ``Victorian Aboriginal languages'': general context. \url{https://museumsvictoria.com.au/bunjilaka/about-us/first-peoples/victorian-aboriginal-languages/}
\item VCAA, Aboriginal Languages of Victoria curriculum history document. \url{https://vcaa.vic.edu.au:443/Documents/alcv/History.pdf}
\end{itemize}

\section{Religious and Community Resources}
None specific to Yabula Yabula were found. Yorta Yorta community language work may be a relevant neighbour, but I did not find sources on it.

\section{Digital and Audio Resources}
No audio or corpus resources were found.

\section{Curriculum Development Context}
The language is extinct and little is documented. A curriculum would draw mainly on Bowe and Morey (1999). It should be developed with the descendant community, and it could be linked to Yorta Yorta revitalization work.
\end{document}